% !PRINTED 196
\weberpage{196}
Dass die Vorzeichen von $\sqrt{-3\,D}$ mit Rücksicht auf die Vorzeichen in (11) so zu nehmen sind, wie es in den Formeln (13) geschehen ist, ergiebt sich am einfachsten, wenn man die Rechnung für den besonderen Fall $a_0 = 0$, $a_3 = 0$ durchführt.

\begin{center}
\S.\,63.
\end{center}

\begin{center}
\gesperrt{Das volle Formensystem der binären cubischen Form.}
\end{center}

Wir wollen noch nachweisen, dass die Invarianten und Covarianten der cubischen Form durch die Functionen $D, f, H, Q$ erschöpft sind, d.\,h. dass alle Invarianten und Covarianten einer cubischen Form sich rational durch diese vier Formen ausdrücken lassen.

Sei also
\[
C = C(x,y,a)
\]
eine Covariante vom Grade $\nu$ in den Variablen und vom Grade $\mu$ in den Co\"efficienten, von der wir voraussetzen wollen, dass sie nur ganzzahlige numerische Co\"efficienten hat.

Es ist dann für irgend eine transformirte Form
\weq{1}{C' = C(\xi,\eta,a') = r^{\lambda}\,C(x,y,a),}
\weq{2}{\lambda = \frac{3\,\mu - \nu}{2}.}

Wir verstehen jetzt unter $\xi$, $\eta$ die Variablen unserer oben betrachteten Normalform. Bezeichnet dann $\varepsilon$ eine dritte Einheitswurzel, so ist
\[
\xi = \varepsilon\,\xi',\quad \eta = \varepsilon^2\,\eta'
\]
eine lineare Substitution mit der Determinante 1, durch die die Normalform \S.\,62, (9) in sich selbst übergeht. Es folgt also, dass $C'$ sich nicht ändern kann, wenn $\xi$, $\eta$ durch $\varepsilon\,\xi$, $\varepsilon^2\,\eta$ ersetzt werden; es kann also $C'$ nur von $\xi\,\eta$, $\xi^3$, $\eta^3$ rational abhängen.

Die Vertauschung von $\xi$ mit $\eta$ entspricht einer Substitution mit der Determinante $-1$; dadurch bleibt also $C'$ ungeändert oder ändert sein Zeichen, je nachdem $\lambda$ gerade oder ungerade ist. Es besteht daher $C'$ aus einer Summe von Gliedern der Form
\[
M\,(\xi\,\eta)^h\,(\xi^{3k} \pm \eta^{3k}),
\]
worin $M$ ein ganzzahliger Factor ist (da wir die numerischen Co\"efficienten in $C$ und also auch in $C'$ als ganze Zahlen voraus-
%CONT>>
